% !TEX root = chapter-standalone.tex

\section{(Co)slice construction}
\label{sec:coslice}

\todotext{@J: review and edit}

The final constructions of this chapter again promote morphisms to objects, but this time we do not take \emph{all} morphisms of a category~\CatC: we fix one object~$\slicetarget$ of special interest, and consider only the morphisms \emph{into} it.

The resulting \emph{slice category} is a general way of modeling ``objects of \CatC equipped with additional information valued in~$\slicetarget$''.
The fixed object~$\slicetarget$ plays the role of a context, a type system, an index, or a base; the objects of the slice are the objects of~\CatC together with their assignment into that context; and the morphisms of the slice are the morphisms of~\CatC that \emph{respect} the assignments.
The examples after the definition make this concrete.

\begin{ctdefinition}[Slice categories]
    \label{def:slice-category}

    Let \CatC be a category and fix an object $\slicetarget \setin \Ob_\CatC$.
    The \maindef{slice category} $\sliceCat{C}{\slicetarget}$ \emph{of} $\CatC$ \emph{over} $\slicetarget$ is:
    \begin{enumerate}
        \item \emph{Objects:} morphisms in \CatC of the type $\Obja \mto \slicetarget$, where $\Obja$ ranges over all objects of $\CatC$.
        \item \emph{Morphisms:} given objects $\Obja \overset{\mora}{\mto} \slicetarget$ and $\Objb \overset{\morb}{\mto} \slicetarget$, a morphism
              \begin{equation}
                  \slicemora_{\sliceCat{C}{\slicetarget}} \colon (\Obja \overset{\mora}{\mto} \slicetarget)
                  \mtoin{\sliceCat{C}{\slicetarget}} (\Objb \overset{\morb}{\mto} \slicetarget)
              \end{equation}
              is specified by a morphism $\slicemora \colon \Obja \mtoin{\CatC} \Objb$ such that the diagram
              \begin{equation}\label{eq:slice-morphism}
                  \includesag{slice-morphism}
              \end{equation}
              commutes.

        \item \emph{Composition:} defined via the composition in $\CatC$.
              Concretely, given composable morphisms $\slicemora_{\sliceCat{C}{\slicetarget}}$ and $\slicemorb_{\sliceCat{C}{\slicetarget}}$ of $\sliceCat{C}{\slicetarget}$, we define
              \begin{equation}
                  \slicemora_{\sliceCat{C}{\slicetarget}} \mthenof{\sliceCat{C}{\slicetarget}} \slicemorb_{\sliceCat{C}{\slicetarget}} \definedas (\slicemora \mthenof{\CatC} \slicemorb)_{\sliceCat{C}{\slicetarget}}.
              \end{equation}
        \item \emph{Identities:} defined by the identities in $\CatC$.
    \end{enumerate}
\end{ctdefinition}

\begin{gradedexercise}[\exname{SliceCat}]
    \label{ex:SliceCat}
    Let~$\CatC$ be a category, fix~$\slicetarget \setin \Ob_{\CatC}$, and consider the slice category~$\sliceCat{\CatC}{\slicetarget}$.
    Your task is to check that the composition of two composable morphisms in~$\sliceCat{\CatC}{\slicetarget}$ is again in fact a morphism in~$\sliceCat{\CatC}{\slicetarget}$.
\end{gradedexercise}

\solutionof{SliceCat}

\begin{example}[Typed sets]
    \label{exa:slice-set-typed}
    Fix a set~$\slicetarget$ and consider the slice category~$\sliceCat{Set}{\slicetarget}$.
    An object is a function~$\mora \colon \Obja \mto \slicetarget$, which we can read as a set of items~$\Obja$ in which every item carries a ``type'' or ``tag'' drawn from~$\slicetarget$.
    A morphism from~$\mora \colon \Obja \mto \slicetarget$ to~$\morb \colon \Objb \mto \slicetarget$ is a function~$\slicemora \colon \Obja \mto \Objb$ with~$\slicemora \mthen \morb = \mora$: it may rename or merge items, but it must preserve each item's type.

    For instance, let~$\slicetarget$ be a set of file formats.
    An object of~$\sliceCat{Set}{\slicetarget}$ is a collection of files, each with its format; a morphism is a reorganization of the files (renaming, deduplication, restructuring) that leaves every file's format unchanged.
    The slice construction packages ``sets fibered over~$\slicetarget$'' as a category in its own right.
\end{example}

\begin{example}[Slices of a poset]
    \label{exa:slice-poset}
    Let~$\posA$ be a \SY{poset}, viewed as a category (\cref{sec:posets-as-cats}), and fix an element~$\slicetarget$.
    Since there is at most one morphism between any two objects, an object of the slice~$\poscat{\posA} / \slicetarget$ is simply an element~$\posela$ with~$\posela \posAleq \slicetarget$, and morphisms are the order relations among such elements.
    The slice is thus the ``down-set'' of~$\slicetarget$: the part of the poset that lies below the chosen element.
    If~$\posA$ orders, say, resource budgets, then slicing over~$\slicetarget$ restricts attention to everything affordable within budget~$\slicetarget$.
\end{example}

\begin{ctdefinition}[Coslice categories]
    \label{def:coslice-category}

    Let \CatC be a category and fix an object~$\coslicesource \setin \Ob_\CatC$.
    The \maindef{coslice category} $\cosliceCat{C}{\coslicesource}$ \emph{of} $\CatC$ \emph{under} $\coslicesource$ is:
    \begin{enumerate}
        \item \emph{Objects:} morphisms in \CatC of the type $\coslicesource \mto \Obja$, where $\Obja$ ranges over all objects of $\CatC$.
        \item \emph{Morphisms:} given objects $\coslicesource \overset{\mora}{\mto} \Obja$ and $\coslicesource  \overset{\morb}{\mto} \Objb$, a morphism
              \begin{equation}
                  \slicemora_{\cosliceCat{C}{\coslicesource}} \colon (\coslicesource \overset{\mora}{\mto} \Obja)
                  \mtoin{\cosliceCat{C}{\coslicesource}}
                  (\coslicesource \overset{\morb}{\mto} \Objb)
              \end{equation}
              is specified by a morphism $\coslicemora \colon \Obja \mtoin{\CatC} \Objb$ such that the diagram
              \begin{equation}\label{eq:coslice-morphism}
                  \includesag{coslice-morphism}
              \end{equation}
              commutes.

        \item \emph{Composition:} defined by the composition in $\CatC$.
        \item \emph{Identities:} defined by the identities in $\CatC$.
    \end{enumerate}
\end{ctdefinition}

